\documentclass[11pt,a4paper]{article}
\usepackage[T1]{fontenc}
\usepackage{lmodern}
\usepackage[margin=27mm]{geometry}
\usepackage{amsmath,amssymb,amsthm,mathtools}
\usepackage{booktabs,array}
\usepackage{microtype}
\usepackage[numbers,sort&compress]{natbib}
\usepackage[hidelinks]{hyperref}
\usepackage{url}
\hypersetup{pdftitle={Four nested limit cycles in a planar power-law reaction system},
pdfauthor={Xiaojian Zeng},pdfsubject={Local draft; human scientific review pending},
pdfkeywords={limit cycles, power-law kinetics, weak focus, interval arithmetic}}
\newcommand{\doi}[1]{\url{https://doi.org/#1}}
\newtheorem{theorem}{Theorem}
\newtheorem{proposition}[theorem]{Proposition}
\newtheorem{lemma}[theorem]{Lemma}
\theoremstyle{remark}
\newtheorem{remark}[theorem]{Remark}
\newcommand{\R}{\mathbb R}
\newcommand{\ii}{\mathrm i}
\newcommand{\Rea}{\operatorname{Re}}
\newcommand{\D}{\mathrm D}
\newcommand{\eps}{\varepsilon}
\newcommand{\bino}[2]{\binom{#1}{#2}}
\numberwithin{equation}{section}
\setlength{\emergencystretch}{2em}
\title{Four nested limit cycles in a planar\texorpdfstring{\\}{ }power-law reaction system}
\author{Xiaojian Zeng\\[3pt]
\small PBC School of Finance, Tsinghua University, 100190\\
\small \texttt{xiaojianzeng@foxmail.com}}
\date{}
\begin{document}
\maketitle
\begin{abstract}
We prove that a planar power-law reaction system considered by Boros and
Hofbauer admits a stable positive equilibrium surrounded by four nested
hyperbolic limit cycles. From the innermost to the outermost, their
stability alternates between repelling and attracting. The system
therefore supports a stable resting state and two stable oscillations at
the same parameter values. The proof starts from a two-parameter
reduction of the Hopf and first-focus conditions. Directed integer
interval arithmetic isolates a common zero of the next two focus
coefficients, proves that the ninth-order radial coefficient is negative,
and verifies a nonsingular four-parameter unfolding. A separate real
Lyapunov-polynomial calculation with directed decimal intervals confirms
the decisive root, sign, and rank conditions. A parameter-dependent finite
normal form and a uniform continuously differentiable remainder estimate
then produce four simple fixed points of the return map of the original
analytic differential equation. All cycles lie in a compact annulus
strictly inside the positive quadrant. The result is existential and
local: it does not specify a finite-amplitude parameter example or a
physical period, and it allows real reaction orders.
\end{abstract}
\noindent\textbf{Keywords:} limit cycles; power-law kinetics; weak focus;
Hopf bifurcation; interval arithmetic.\\
\textbf{Mathematics Subject Classification:} 34C07, 34C23, 80A30.

\section{Introduction}
Uniqueness of a positive steady state does not preclude several attracting
oscillations in a nonlinear reaction model. A particularly economical
example was proposed by Boros and Hofbauer \cite[Section~5.1]{BorosHofbauer2024}.
They identified approximate parameters at which three focus values vanish
and a fourth is negative, and left the existence of four limit cycles
as a conjecture. We establish the corresponding existence statement for
their specified four-parameter family.

The issue is more than the accuracy of a numerical weak-focus location.
One must prove that the simultaneous zero exists, certify the sign of the
next coefficient, verify enough independent parameter directions, and
control the difference between a finite normal form and the original
equation. We supply these four ingredients. The finite coefficient
calculations are checked in two different representations and two
different arithmetic implementations; the final passage to periodic
orbits uses the true return map.

Multiple nested reaction-system cycles, including alternating stability,
are already known in other families. In particular, Craciun and Erban
\cite[Theorem~1, equations~(4.27)--(4.30), and Section~5]{CraciunErban2025}
construct nested algebraic cycles using polynomial products. Our
contribution concerns the particular sparse power-law family below,
rather than the general coexistence pattern or the general method of
unfolding a weak focus. The calculation does not depend on the optional
notebook cited by Boros and Hofbauer, whose contents were unavailable for
comparison. No exhaustive priority claim about that material is made.

\section{Model and result}\label{sec:model}
Consider, in the original dimensionless time $t$,
\begin{equation}\label{eq:ode}
 \dot x=y^{-1}-x^a y^b,\qquad
 \dot y=K\bigl[-(d-1)-y^{-1}+d x^a y^b\bigr]
\end{equation}
on $\R_{>0}^2$, with parameter domain
\begin{equation}\label{eq:domain}
 \mathcal P=\{(K,a,b,d):K>0,\ a>0,\ b>-1,\ d>1,\ 1+bd>0\}.
\end{equation}
These are the normalized equations and inequalities of
\cite[Section~5.1]{BorosHofbauer2024}, with $d>1$ enforcing existence of
the positive steady state. Real powers have their usual positive-real
meaning. Consequently the vector field is jointly real analytic in
state and parameters on its stated domain. No boundary condition on an
axis is imposed or needed for the local construction.

Solving the first equilibrium equation gives $x^a y^b=y^{-1}$; the second
then gives $(d-1)(y^{-1}-1)=0$. Thus the unique positive equilibrium is
$(1,1)$. Its Jacobian, trace, and determinant are
\begin{equation}\label{eq:linear}
 J=\begin{pmatrix}-a&-(1+b)\\Kda&K(1+bd)\end{pmatrix},\qquad
 T=-a+K(1+bd),\qquad \Delta=Ka(d-1)>0.
\end{equation}

\begin{theorem}\label{thm:main}
There exist parameters in $\mathcal P$ for which the unique positive
equilibrium of \eqref{eq:ode} is hyperbolically asymptotically stable
and is surrounded by at least four disjoint nested hyperbolic periodic
orbits. From the innermost to the outermost, the four orbits are repelling,
attracting, repelling, and attracting. They are all contained in a compact
annulus in $\R_{>0}^2$ that has positive distance from $(1,1)$ and from
the coordinate axes.
\end{theorem}

The parameters will approach a rigorously isolated fourth-order weak
focus $p_*$, approximately
\begin{equation}\label{eq:pstar}
 \begin{split}
 a_*&\simeq1.0128219986803146,\qquad b_*\simeq0.6546256211042382,\\
 d_*&\simeq3.2886104393099358,\qquad K_*\simeq0.3212443603989608.
 \end{split}
\end{equation}
These rounded values locate the organizing weak focus, not a parameter
example with four cycles. The cycles occur at nearby parameters
constructed in Section~\ref{sec:cycles}. Hyperbolicity refers to the
nontrivial multiplier of each periodic orbit.

\section{Focus coefficients and reduction}\label{sec:coefficients}
Write $q=1+bd$. On the Hopf hypersurface $K=a/q$, put
$\omega^2=\Delta$ and set
\begin{equation}\label{eq:basis}
 u=x-1,\quad v=y-1,\quad X=u,\quad
 Y=\frac{au+(1+b)v}{\omega},\quad z=X+\ii Y.
\end{equation}
The linear field is $\dot z=\ii\omega z$ with positive angular orientation.
All coefficient definitions below retain time $t$.

For $i+j\geq1$, the coefficients of the original real Taylor jet are
\begin{equation}\label{eq:jet}
 \begin{split}
 f_{ij}&=\mathbf1_{i=0}\bino{-1}{j}-\bino a i\bino b j,\\
 g_{ij}&=K\left[-\mathbf1_{i=0}\bino{-1}{j}
                  +d\bino a i\bino b j\right].
 \end{split}
\end{equation}
Here $\binom\alpha0=1$ and
$\binom\alpha{j+1}=\binom\alpha j(\alpha-j)/(j+1)$; the constant
coefficients cancel. Setting $s=\ii\omega$, substitute
\begin{equation}\label{eq:complex-jet}
 u=\frac{z+\bar z}{2},\qquad
 v=\frac{(-a-s)z+(-a+s)\bar z}{2(1+b)},\qquad
 \mathcal F=f+\frac{s}{\omega^2}\bigl[af+(1+b)g\bigr].
\end{equation}
Computations in the algebra $s^2=-\omega^2$ use rational operations only.

We fix the following finite normal-form convention. Find $H=z+O(|z|^2)$
and $G=sz+O(|z|^2)$ satisfying, through degree nine,
\begin{equation}\label{eq:conjugacy}
 \mathcal F(H,\bar H)=H_zG+H_{\bar z}\bar G.
\end{equation}
All resonant coefficients $H_{m+1,m}$ are set to zero. If degrees below
$n$ have been determined, let $\mathcal R_n$ be the degree-$n$ residual
of the left side minus the right side of \eqref{eq:conjugacy}. For
$j+k=n$, the next coefficients are
\begin{equation}\label{eq:recurrence}
 \begin{cases}
 G_{jk}=(\mathcal R_n)_{jk},\quad H_{jk}=0,&j-k=1,\\[2pt]
 H_{jk}=\dfrac{(\mathcal R_n)_{jk}}{(j-k-1)s},\quad G_{jk}=0,&j-k\ne1.
 \end{cases}
\end{equation}
Indeed, adding $H_{jk}$ changes this residual by
$-(j-k-1)sH_{jk}$, and adding $G_{jk}$ changes it by $-G_{jk}$.
Every nonresonant denominator is nonzero. Define
\begin{equation}\label{eq:focus}
 F_m=\Rea G_{m+1,m},\qquad 1\leq m\leq4.
\end{equation}
If $F_1,\ldots,F_{m-1}$ vanish, the first nonzero radial return
coefficient is $(2\pi/\omega)F_m$. This follows by writing
$z=r e^{\ii\theta}$ and integrating the first radial drift over one
positive angular revolution. Coefficients beyond the first nonzero one
need not agree across different normal-form conventions.

The degree-three instance of \eqref{eq:jet}--\eqref{eq:recurrence} gives
the rational identity
\begin{equation}\label{eq:F1}
 F_1=\frac{a(2a^2d+aq-q^2)}{16q^2}.
\end{equation}
After multiplying by $2\pi/\omega$, this is the first focus value
printed in \cite[Section~5.1]{BorosHofbauer2024}. Introduce $h>2$ by
\begin{equation}\label{eq:reduction}
 d=\frac{h(h-1)}2,\qquad b=\frac{ah-1}{d},\qquad K=\frac1h.
\end{equation}
Then $q=ah$, so $T=F_1=0$ identically. Taking the derivative with respect
to $b$ on the Hopf hypersurface first, and then imposing
\eqref{eq:reduction}, yields
\begin{equation}\label{eq:bderivative}
 \left.\partial_bF_1\right|_{T=0,\,\eqref{eq:reduction}}
       =-\frac{(h-1)(2h-1)}{32h}<0.
\end{equation}
Thus eliminating the first focus equation has not removed its transverse
parameter direction. In what follows $f(a,h)=(F_2,F_3)$ denotes the
restriction to \eqref{eq:reduction}.

\subsection{Exact calibrations}
The calculation has two useful published checks. For the system of
\cite[equation~(12), Proposition~3]{BorosHofbauer2024},
\begin{equation}\label{eq:calibration}
 \begin{split}
 \dot x&=y^{-q_c}-xy^{1/2},\\
 \dot y&=K_c\bigl[-(q_c+r_c+1/2)+(q_c+1/2)y^{-q_c}+r_cxy^{1/2}\bigr],
 \end{split}
\end{equation}
at $(q_c,r_c,K_c)=(1/4,15/8,4/3)$, the convention above gives
\begin{equation}\label{eq:benchmark}
 \omega^2=\frac78,\qquad F_1=F_2=0,\qquad
 F_3=-\frac{4375}{884736},\qquad
 \frac{2\pi F_3}{\omega}=-\frac{625\pi}{110592}\sqrt{\frac72}.
\end{equation}
For the endpoint $(a,h)=(1,22/7)$ of the target family, corresponding
to $d=165/49$, exact rational arithmetic instead gives
\begin{equation}\label{eq:endpoint}
 F_1=F_2=0,\qquad F_3=-\frac{18125}{2404233216}<0.
\end{equation}
The zero and sign in \eqref{eq:endpoint} agree with the boundary
statement in \cite[Section~5.1]{BorosHofbauer2024}; its displayed fraction
is our calculation in convention \eqref{eq:focus}. The two benchmark
systems are distinct.

\section{An interval certificate for the weak focus}\label{sec:certificate}
The following two finite decimals are exact rational inputs:
\begin{equation}\label{eq:center}
 \begin{split}
 a_0&=1.012821998680314600068208268899051820974352006,\\
 h_0&=3.112895114354931694231478861591838762886663535.
 \end{split}
\end{equation}
Let $x_0=(a_0,h_0)$ and $\delta=10^{-20}$. For the interval calculation
we use an outward dyadic enclosure $\mathcal X$ of
$x_0+[-\delta,\delta]^2$, with grid spacing $2^{-384}$. The exact
rational preconditioner is given entrywise by
\begin{equation}\label{eq:preconditioner}
 \begin{split}
 R_{11}&=-523.0513662343578025786281282452785879709,\\
 R_{12}&= 1724.438413838856510805394172256443795935,\\
 R_{21}&= 1632.401201395760086859367561999686168763,\\
 R_{22}&=-3959.57512664182068880590301715606501281.
 \end{split}
\end{equation}
Its determinant is nonzero by exact rational arithmetic. Neither the
center nor the preconditioner is assumed accurate in advance.

\begin{proposition}\label{prop:certificate}
The box $x_0+[-10^{-20},10^{-20}]^2$ contains a unique common zero
$(a_*,h_*)$ of $(F_2,F_3)$. At this zero $F_4<0$, and
$\det\D_{(a,h)}(F_2,F_3)\ne0$. The corresponding parameters
\eqref{eq:reduction} belong to $\mathcal P$.
\end{proposition}

\begin{proof}
Evaluate \eqref{eq:jet}--\eqref{eq:recurrence} with integer interval
endpoints on the $2^{-384}$ grid. Addition and subtraction are exact on
this grid. For multiplication and division, all endpoint products or
quotients are considered and rounded outward using integer floor and
ceiling. A division is allowed only when the denominator interval
excludes zero. All divisions in this calculation satisfy that check.
First derivatives with respect to $a,h$ are propagated by the product
and quotient rules along the same recurrence. The resulting interval
matrix $J_f(\mathcal X)$ therefore encloses $\D f(x)$ at every
$x\in\mathcal X$.

Table~\ref{tab:certificate} records convenient rational outer bounds
on the resulting intervals. Full integer endpoints and the executable
recurrence are supplied in Online Resource~1. In particular, the table
is a rounded outward display of strict enclosures, not floating-point
error estimates.
\begin{table}[htbp]
\centering
\caption{Directed dyadic enclosures on $\mathcal X$. Every displayed
interval has been widened from the exact computed endpoints.}
\label{tab:certificate}
\small
\begin{tabular}{@{}ll@{}}
\toprule
Quantity & Rational outer bound\\
\midrule
$(J_f)_{11}$ & $[0.005322623692367449,\ 0.005322623692367584]$\\
$(J_f)_{12}$ & $[0.002318061020176081,\ 0.002318061020176173]$\\
$(J_f)_{21}$ & $[0.002194340814884875,\ 0.002194340814887489]$\\
$(J_f)_{22}$ & $[0.000703107152964103,\ 0.000703107152965913]$\\
$\det J_f$ & $[-1.344241117333,\ -1.344241117316]\,10^{-6}$\\
$F_4$ & $[-3.54349145,\ -3.54349141]\,10^{-7}$\\
First row sum of $|I-RJ_f|$ & $<3.872\,10^{-12}$\\
Second row sum of $|I-RJ_f|$ & $<8.938\,10^{-12}$\\
Both coordinates of $|\mathcal Y-x_0|$ & $<9\,10^{-32}$\\
\bottomrule
\end{tabular}
\end{table}

Here the interval image is
\begin{equation}\label{eq:image}
 \mathcal Y=x_0-Rf(x_0)
                  +(I-RJ_f(\mathcal X))(\mathcal X-x_0).
\end{equation}
The map $\Phi(x)=x-Rf(x)$ maps $\mathcal X$ into $\mathcal Y$ by
the integral mean-value formula. The row-sum bounds make it a strict
contraction in the maximum norm, and the image bound places
$\mathcal Y$ strictly inside the original rational box. Banach's
fixed-point theorem gives one and only one fixed point in $\mathcal X$.
Since $R$ is nonsingular, this is exactly a zero of $f$. The fixed point
belongs to $\mathcal Y$, so it also lies in the original box without
dyadic enlargement.

The determinant and $F_4$ intervals prove the two remaining coefficient
claims. Direct interval evaluation of \eqref{eq:reduction} gives
$a>0$, $b>-1$, $d>1$, $K>0$, $1+bd>0$, and $\omega^2>0$ on the
whole box. For example, deliberately coarse enclosures are
$1<a<1.02$, $0.65<b<0.66$, $3.28<d<3.30$, $0.32<K<0.33$,
and $0.74<\omega^2<0.75$. This also establishes admissibility.
\end{proof}

The code uses only the Python standard library for strict certification.
Its input is \eqref{eq:center}--\eqref{eq:preconditioner}, not previously
calculated focus values. A separate real scalar recurrence and directed
decimal interval implementation certify the same zero and signs;
Appendix~\ref{app:independent} states its conventions and bounds.

\section{A full parameter unfolding}\label{sec:unfolding}
First extend the finite normal form smoothly off the Hopf hypersurface.
For $p=(K,a,b,d)$ near $p_*$, set
\begin{equation}\label{eq:offhopf}
 \tau=T/2,\qquad \nu=\sqrt{\Delta-\tau^2}>0,\qquad
 X=u,\quad Y=\frac{(a+\tau)u+(1+b)v}{\nu}.
\end{equation}
The linear field is $(\tau X-\nu Y,\nu X+\tau Y)$.
With $\lambda=\tau+\ii\nu$, use
$j\lambda+k\bar\lambda-\lambda$ as the nonresonant denominator in
\eqref{eq:recurrence}, still retaining all terms with $j-k=1$.
For $j-k\ne1$, the imaginary part $(j-k-1)\nu$ is bounded away
from zero after restricting the parameter neighborhood. This is a
finite polynomial change of variables depending smoothly on $p$.
Its linear part is the identity, so it is a local diffeomorphism on
a uniform sufficiently small state neighborhood.

The actual transformed field, including its remainder, has the form
\begin{equation}\label{eq:fullnf}
 \dot z=(\tau+\ii\nu)z+\sum_{m=1}^4 C_m(p)z|z|^{2m}+E_p(z),
\end{equation}
where $\Rea C_m=F_m$ on $T=0$. On a sufficiently small closed
parameter neighborhood and state ball there is a uniform constant $C$
such that
\begin{equation}\label{eq:remainder}
 |E_p(z)|\leq C|z|^{10},\qquad
 \|\D_zE_p(z)\|\leq C|z|^9.
\end{equation}
These estimates follow from Taylor's theorem for the jointly analytic
original field and the finite coordinate change on a compact set.
They do not assume convergence of an infinite normal-form series.

Define
\begin{equation}\label{eq:beta}
 \beta(p)=(\beta_0,\beta_1,\beta_2,\beta_3)
      =(T/2,\Rea C_1,\Rea C_2,\Rea C_3).
\end{equation}
We next verify that these four coefficients are local parameter
coordinates. Replace $(K,a,b,d)$ by $(T,b,a,h)$, where
$K=(T+a)/(1+bd)$ and $d=h(h-1)/2$. On $T=0$, the first-focus
zero is the graph in \eqref{eq:reduction}. By the chain rule, or
equivalently the Schur complement,
\begin{equation}\label{eq:schur}
 \det\frac{\partial(F_1,F_2,F_3)}{\partial(b,a,h)}
   = (\partial_bF_1)\det\D_{(a,h)} f
\end{equation}
at $p_*$. Accounting for $\tau=T/2$ and the parameter order gives
\begin{equation}\label{eq:fullrank}
 \det\D_{(K,a,b,d)}\beta(p_*)
 =-\frac{1+bd}{2h-1}(\partial_bF_1)\det\D f(a_*,h_*)\ne0.
\end{equation}
The factors are nonzero by \eqref{eq:bderivative} and
Proposition~\ref{prop:certificate}; a strict outer enclosure of the
determinant is $[-8.990,-8.989]\,10^{-8}$. Derivatives of $C_m$ in
the $T$ direction do not change this triangular determinant argument.
The inverse function theorem therefore applies to $\beta$ at $p_*$.

The unfolding varies $a,b,d$ as well as $K$: reaction powers and a
reaction-vector component change. It is not an unfolding by rate
constants of a single fixed network. Thus it does not contradict the
rate-constant codimension bound in
\cite[Theorem~35]{BanajiBorosHofbauer2024}. Likewise, the inheritance
theorem in \cite[Theorem~3.2]{BanajiBorosHofbauer2025} assumes an already
transverse source bifurcation; it does not replace the coefficient and
rank calculation above.

\section{Four cycles of the original analytic equation}\label{sec:cycles}
\begin{proof}[Proof of Theorem~\ref{thm:main}]
Let $\ell=F_4(p_*)<0$ and $\omega_*=\sqrt{\Delta(p_*)}>0$.
By \eqref{eq:fullrank}, for every sufficiently small $\eps>0$ there
is a fixed parameter vector $p(\eps)$ satisfying
\begin{equation}\label{eq:unfolding}
 \begin{aligned}
 \beta_0&=24\ell\eps^8,&\qquad \beta_1&=-50\ell\eps^6,\\
 \beta_2&=35\ell\eps^4,&\qquad \beta_3&=-10\ell\eps^2.
 \end{aligned}
\end{equation}
Moreover $p(\eps)-p_*=O(\eps^2)$, so
$\Rea C_4(p(\eps))=\ell+O(\eps^2)$. These parameters remain in
the open set $\mathcal P$. The equilibrium has eigenvalues
$\beta_0\pm\ii\nu$ with $\beta_0<0$, hence is hyperbolically stable.

Put $z=\eps R e^{\ii\theta}$ and consider a fixed compact $R$
interval away from zero, such as $[1/2,5/2]$, together with a slightly
larger interval for controlling the flow. Equations
\eqref{eq:fullnf}--\eqref{eq:unfolding} imply
$\dot\theta=\omega_*+O(\eps^2)>0$. Define
\begin{equation}\label{eq:polynomial}
 P(R)=R(R^2-1)(R^2-2)(R^2-3)(R^2-4).
\end{equation}
The identity
$\prod_{j=1}^4(S-j)=S^4-10S^3+35S^2-50S+24$ and the uniform
remainder bounds give
\begin{equation}\label{eq:radial}
 \frac{dR}{d\theta}
       =\eps^8\frac{\ell}{\omega_*}P(R)+O_{C^1}(\eps^9).
\end{equation}
The notation $O_{C^1}$ means a uniform bound on the function and
its first $R$ derivative, also uniform in $\theta$. In particular,
the radial projection of $E_p$ is $O(\eps^{10})$; dividing by
$\eps$ makes its contribution $O(\eps^9)$, with the same order
for the $R$ derivative by \eqref{eq:remainder}. The angular-speed
correction and the change in $C_4$ contribute only $O(\eps^{10})$
to \eqref{eq:radial}.

Let $\Pi_\eps$ be the actual return map on $\theta=0$, after an
increase of $2\pi$ in $\theta$. Its definition uses solutions of
\eqref{eq:fullnf} including $E_p$, hence solutions of
\eqref{eq:ode}. Over this fixed angular interval, $R$ changes by
$O(\eps^8)$. Integrating \eqref{eq:radial} and its variational
equation gives
\begin{equation}\label{eq:return}
 \frac{\Pi_\eps(R)-R}{\eps^8}
       =\frac{2\pi\ell}{\omega_*}P(R)+O_{C^1}(\eps).
\end{equation}
For completeness, a uniform bound on the right side of
\eqref{eq:radial} keeps the flow in the slightly larger interval
through $2\pi$; its $R$ derivative is $O(\eps^8)$, so the
variational solution is $1+O(\eps^8)$ by Gronwall's inequality.
Substitution into the integral equations yields both estimates in
\eqref{eq:return}. No long-time numerical approximation is involved.

The four positive roots $R_j=\sqrt j$ of $P$ are simple. Choose
four disjoint small closed intervals around them. On each interval
the limiting displacement in \eqref{eq:return} has opposite endpoint
signs and a derivative of constant nonzero sign. Uniform $C^1$
convergence preserves these properties for sufficiently small fixed
$\eps>0$. There is exactly one simple return-map fixed point in
each of the four intervals.

The derivatives $P'(R_j)$ are $-12,8,-12,48$ in order. Since
$\ell<0$, the signs of $\Pi_\eps'-1$ at the four fixed points are
$+,-,+,-$. Each multiplier is positive and approaches one as
$\eps\to0$, so these signs give the asserted alternating
repelling and attracting stability. Positive angular speed makes
each periodic orbit wind once around the equilibrium. Uniqueness
of planar solutions makes the four Jordan curves disjoint and nested.
The local coordinate transformations preserve this geometry and
orbital stability.

Finally fix one sufficiently small positive $\eps$. All four curves
lie in a neighborhood of $(1,1)$ with closure in the positive
quadrant. Their compact union avoids the equilibrium and the axes.
Choosing an inner Jordan curve inside the innermost orbit and an
outer Jordan curve outside the outermost orbit gives the required
compact annulus. This proves the theorem for the original equation.
\end{proof}

\subsection{Physical meaning and limitations}
At the constructed parameters, a stable rest state and two attracting
oscillations coexist in the same dimensionless two-concentration
model. The two repelling cycles are local separators of the nearby
attracting dynamics. The result concerns a neighborhood of the steady
state; it gives neither global permanence nor an upper bound on the
total number of cycles in the positive quadrant.

The finite normal form determines a sufficiently small parameter
unfolding, while analyticity supplies the existence of a uniform
remainder constant. We have not evaluated that constant to give a
numerical upper bound for $\eps$. Thus \eqref{eq:pstar} must not be
used as a directly simulated four-cycle example. Computing explicit
finite-amplitude parameters and enclosing their four return roots is
a further task. No physical period is reported. The argument retains
the dimensionless time in \eqref{eq:ode}; a positive monomial change
of time would preserve orbits but would not preserve periods.
Finally, real reaction orders and the permitted reaction-vector
variation do not establish an integer-stoichiometric chemical
implementation or an experimentally observed oscillator.

\appendix
\section{Independent scalar Lyapunov calculation}\label{app:independent}
A second calculation avoids the complex conjugacy recurrence. On the
Hopf hypersurface use $X=u$, $Z=au+(1+b)v$, $w=\omega^2$, and
$Q=wX^2+Z^2$. The linear field is $(-Z,wX)$. Reconstruct the real
field directly from \eqref{eq:jet} and $v=(Z-aX)/(1+b)$ and seek
\begin{equation}\label{eq:lyapunov}
 H=Q+\sum_{n=3}^{10}H_n,\qquad
 \nabla H\cdot(\dot X,\dot Z)=\sum_{m=2}^5 c_m Q^m
                                    +O(\|(X,Z)\|^{11}),
\end{equation}
where $H_n$ is homogeneous of degree $n$. Although $H$ is computed
through degree ten, only the degree-nine vector-field jet is needed.

Let $B_n$ be the contribution at degree $n$ of the previously computed
$H_k$ and the nonlinear field, and let
$L=-Z\partial_X+wX\partial_Z$. For odd $n$ solve $LH_n=-B_n$.
For $n=2m$, set $c_m$ equal to the angular mean of $B_n$ on $Q=1$
and solve $LH_n=c_mQ^m-B_n$. The required averages are
\begin{equation}\label{eq:average}
 \left\langle X^{2p}Z^{2q}\right\rangle_{Q=1}
 =\frac{(2p)!(2q)!}{4^{p+q}p!q!(p+q)!w^p},
\end{equation}
and are zero when either exponent is odd. If
$H_n=\sum_{j=0}^n h_jX^{n-j}Z^j$, the scalar recurrence is
\begin{equation}\label{eq:scalar}
 w(j+1)h_{j+1}-(n-j+1)h_{j-1}=(\mathrm{rhs})_j,
\end{equation}
with coefficients outside $0\leq j\leq n$ equal to zero. At even
degree impose $h_0=0$ to fix the radial kernel; at odd degree the
last equation fixes $h_0$. These equations give an independent
finite coefficient generator.

The scalar and complex coefficients have a triangular relation. If
$F_1=\cdots=F_{k-1}=0$, the first nonzero term satisfies
\begin{equation}\label{eq:conversion}
 c_{k+1}=\frac{2F_k}{w^k}.
\end{equation}
To see this, pull $H$ into the finite complex normal form. Its
quadratic term is $wr^2$, whose first radial derivative term is
$2wF_kr^{2k+2}$. Division by $Q^{k+1}=w^{k+1}r^{2k+2}$ gives
\eqref{eq:conversion}. Changes in lower radial kernel terms add
only multiples of lower focus values. This proves a smooth
invertible lower-triangular change between the two focus sequences.
In particular, on their common lower-zero locus,
\begin{equation}\label{eq:detconversion}
 \det\D(c_3,c_4)=4w^{-5}\det\D(F_2,F_3),\qquad
 F_4=\tfrac12w^4c_5.
\end{equation}
The derivative determinants here use the reduced variables $(a,h)$.
These equalities are used at the simultaneous zero; a general
off-zero high-order coefficient is not identified across gauges.

The second implementation uses 90-digit decimal endpoints, with
floor and ceiling rounding at every interval operation. It computes
its own exact rational preconditioner by inverting the rational
midpoint of its center derivative enclosure. Its only shared
numerical proposals are \eqref{eq:center} and the radius $10^{-20}$.
Table~\ref{tab:independent} summarizes its strict outer bounds.
\begin{table}[htbp]
\centering
\caption{Independent directed-decimal calculation. The last two
rows use \eqref{eq:detconversion} only at the certified common zero.}
\label{tab:independent}
\small
\begin{tabular}{@{}ll@{}}
\toprule
Quantity & Rational outer bound\\
\midrule
Two contraction row sums & $<2.802\,10^{-11}$, $<6.423\,10^{-11}$\\
Two image displacements & $<2.802\,10^{-31}$, $<6.423\,10^{-31}$\\
$\det\D(c_3,c_4)$ & $[-2.3487381622,-2.3487381619]\,10^{-5}$\\
$c_5$ & $[-2.305150908,-2.305150598]\,10^{-6}$\\
$F_4$ at the common zero & $[-3.543491671,-3.543491195]\,10^{-7}$\\
$\det\D(F_2,F_3)$ at that zero & $[-1.344241118,-1.344241117]\,10^{-6}$\\
\bottomrule
\end{tabular}
\end{table}

The same contraction argument establishes the zero independently.
Exact rational evaluations also give $c_2=c_3=0$ and
$c_4=-625/42336$ for \eqref{eq:calibration}, reproducing
\eqref{eq:benchmark}, and reproduce \eqref{eq:endpoint}.
All homogeneous-equation residuals in those rational checks vanish
exactly. The independent implementation does not import the complex
coefficient generator, its arithmetic module, its preconditioner,
or cached focus outputs.

\section*{Statements and declarations}
\paragraph{Data and code availability.}
Online Resource~1 contains the two coefficient generators, directed
arithmetic implementations, exact rational certificate inputs,
reference output bounds, and instructions for reproducing the
certificates and exact calibrations. No experimental data are used.
No third-party paper is included in the supplementary package.

\paragraph{Use of artificial intelligence.}
OpenAI Codex was used substantively for literature analysis, candidate
derivations, implementation of the calculations, computational checks,
and drafting. A separate AI session implemented the scalar verification
described in Appendix~\ref{app:independent}. These automated checks
are not external peer review and do not establish completed human
scientific review. Human review of the reasoning, code, citations, and
final text remains pending for this local draft.

\paragraph{Author declarations.}
Final author confirmation of affiliation, contributions, funding,
competing interests, and approval of the manuscript is pending.

\bibliographystyle{plainnat}
\bibliography{references}
\end{document}
